\documentclass[border=10pt,tikz]{standalone}
\usepackage{fontspec}
\setmainfont{Apple SD Gothic Neo}
\setsansfont{Apple SD Gothic Neo}
\usepackage{tikz}
\usetikzlibrary{positioning,arrows.meta,shapes,fit,backgrounds}
\begin{document}
% _design.tex — shared TikZ design system for all blog diagrams
% Inputs nothing on its own — meant to be \input{} from each diagram .tex
% AFTER \begin{document}, BEFORE \begin{tikzpicture}.
%
% Provides a unified visual language across GoF / DSA / EMC++ / Beautiful series:
%   - Colors: muted, accessible, color-coded by role
%   - Node styles: consistent sizing, rounded, subtle borders
%   - Edge styles: UML-style inheritance, aggregation, plus dashed alternates
%
% Usage:
%   \documentclass[border=10pt,tikz]{standalone}
%   \usepackage{tikz}
%   \usetikzlibrary{positioning,arrows.meta,shapes,fit,backgrounds}
%   \begin{document}
%   \input{../_design.tex}     % adjust relative path
%   \begin{tikzpicture}[blog]  % apply 'blog' style
%   ...

\definecolor{absfill}{RGB}{220,232,247}   % cool blue
\definecolor{absbord}{RGB}{72,118,182}
\definecolor{confill}{RGB}{249,222,222}   % warm red
\definecolor{conbord}{RGB}{197,93,93}
\definecolor{impfill}{RGB}{223,238,222}   % soft green
\definecolor{impbord}{RGB}{99,154,93}
\definecolor{clifill}{RGB}{251,239,205}   % pale yellow
\definecolor{clibord}{RGB}{200,167,72}
\definecolor{hifill} {RGB}{253,224,200}   % accent orange
\definecolor{hibord} {RGB}{210,135,68}
\definecolor{nodefill}{RGB}{250,250,250}
\definecolor{nodebord}{RGB}{120,120,120}
\definecolor{edgec}  {RGB}{80,80,80}
\definecolor{edgesc} {RGB}{145,145,145}
\definecolor{groupbord}{RGB}{200,200,200}

\tikzset{
  blog/.style={
    font=\sffamily\small,
    every node/.style={font=\sffamily\small}
  },
  % --- node roles ---
  cls/.style={
    rectangle, rounded corners=2pt,
    draw=nodebord, line width=0.4pt,
    minimum height=8mm, minimum width=28mm,
    inner sep=3pt, fill=nodefill, align=center
  },
  abs/.style={cls, draw=absbord, fill=absfill, font=\sffamily\itshape\small},
  con/.style={cls, draw=conbord, fill=confill},
  imp/.style={cls, draw=impbord, fill=impfill},
  cli/.style={cls, draw=clibord, fill=clifill},
  hi/.style={cls, draw=hibord, fill=hifill, font=\sffamily\bfseries\small},
  % --- circle (for trees / graph nodes) ---
  vx/.style={
    circle, draw=absbord, line width=0.4pt,
    minimum size=8mm, fill=absfill, inner sep=0pt
  },
  vxc/.style={vx, draw=conbord, fill=confill},
  vxh/.style={vx, draw=hibord, fill=hifill, font=\sffamily\bfseries},
  % --- decision diamond ---
  q/.style={
    diamond, draw=clibord, line width=0.4pt, fill=clifill,
    inner sep=2pt, aspect=2.4, align=center
  },
  % --- edges ---
  inh/.style={                                  % UML inheritance
    -{Triangle[length=3mm,width=3mm,open]},
    draw=edgec, line width=0.5pt
  },
  agg/.style={                                  % UML aggregation
    {Diamond[length=2.5mm,width=2.5mm,open]}-{Stealth[length=2.2mm]},
    draw=edgec, line width=0.45pt
  },
  uses/.style={                                 % directed reference
    -{Stealth[length=2.2mm]}, draw=edgec, line width=0.5pt
  },
  arr/.style={                                  % plain arrow (graphs/trees)
    -{Stealth[length=2mm]}, draw=edgec, line width=0.45pt
  },
  edge/.style={                                 % undirected (trees)
    draw=edgec, line width=0.45pt
  },
  alt/.style={                                  % dashed: similar/alt
    -{Stealth[length=2mm]}, dashed, draw=edgesc, line width=0.45pt
  },
  ret/.style={                                  % dashed creates/returns
    -{Stealth[length=2mm]}, dashed, draw=edgesc, line width=0.45pt
  },
  thr/.style={                                  % threading link (red dashed)
    ->, dashed, draw=conbord, line width=0.5pt
  },
  % --- group / region ---
  group/.style={
    rectangle, rounded corners=4pt,
    draw=groupbord, line width=0.4pt,
    inner sep=10pt, fill=black!2
  },
  glabel/.style={font=\sffamily\bfseries\small, anchor=north west, inner sep=2pt}
}

% Korean font convenience macro — included by every Hangul diagram
\providecommand{\KOR}{}

\begin{tikzpicture}[blog,
  tc/.style={cls, minimum width=26mm, minimum height=8mm, font=\sffamily\footnotesize},
  mac/.style={cls, minimum width=30mm, minimum height=11mm, align=center},
  seg/.style={rectangle, draw, line width=0.4pt, minimum height=10mm, anchor=west, inner sep=2pt, font=\sffamily\footnotesize, align=center}
]

\node[font=\sffamily\bfseries\large, text=inkc] at (6.0,4.6) {802.1Qbu · 802.3br --- express와 preemptible};

% ---------- upper: structure ----------
\node[font=\sffamily\bfseries\small, text=inkc, anchor=west] at (-1.6,3.6) {(a) 송신 경로};

\node[tc, draw=hibord, fill=hifill] (tce) at (0,2.4) {express TC};
\node[tc, draw=absbord, fill=absfill] (tcp) at (0,0.8) {preemptible TC};

\node[mac, draw=hibord, fill=hifill] (emac) at (4.4,2.4) {eMAC\\[2pt]\footnotesize express MAC};
\node[mac, draw=absbord, fill=absfill] (pmac) at (4.4,0.8) {pMAC\\[2pt]\footnotesize preemptible MAC};

\node[cls, draw=prpbord, fill=prpfill, minimum width=30mm, minimum height=27mm, align=center] (mm) at (8.6,1.6)
  {MAC Merge\\[2pt]\footnotesize (802.3br)};

\node[neut, minimum width=18mm, minimum height=11mm] (phy) at (12.0,1.6) {PHY\\[2pt]\footnotesize 링크};

\draw[arr] (tce.east) -- (emac.west);
\draw[arr] (tcp.east) -- (pmac.west);
\draw[arr] (emac.east) -- (emac.east -| mm.west);
\draw[arr] (pmac.east) -- (pmac.east -| mm.west);
\draw[arr] (mm.east) -- (phy.west);

\node[caption, anchor=north] at (2.2,0.0) {TC마다 E/P 지정 (802.1Q-2018 6.7.2)};

% ---------- lower: timeline ----------
\node[font=\sffamily\bfseries\small, text=inkc, anchor=west] at (-1.6,-0.8) {(b) 링크 위의 개념 timeline};

\node[seg, draw=absbord, fill=absfill, minimum width=36mm] (f1) at (0,-2.6)
  {preemptible 프레임\\[2pt]non-final fragment};
\node[seg, draw=hibord, fill=hifill, minimum width=30mm] (ex) at (3.9,-2.6)
  {express 프레임};
\node[seg, draw=absbord, fill=absfill, minimum width=40mm] (f2) at (7.2,-2.6)
  {preemptible 프레임\\[2pt]남은 fragment};

% express arrival marker
\draw[arr, draw=hibord, line width=0.7pt] (2.2,-1.6) -- (2.2,-2.05);
\node[font=\sffamily\footnotesize, text=hibord, anchor=south] at (2.2,-1.65) {express 프레임 도착};

% residual blocking brace-like marker
\draw[draw=focusc, line width=0.6pt, |-|] (2.2,-3.45) -- (3.9,-3.45);
\node[font=\sffamily\footnotesize, text=focusc, anchor=north, align=center] at (3.05,-3.6)
  {잔여 blocking\\[2pt]속도·fragment 조건·구현에 따름};

% time axis
\draw[arr, line width=0.5pt] (0,-4.9) -- (12.0,-4.9);
\node[caption, anchor=west] at (12.1,-4.9) {시간};

\node[caption, anchor=north, align=center] at (6.0,-5.2)
  {fragment 크기와 지연 값은 링크 양 끝의 지원·설정과 링크 속도로 정해지므로 그림에 수치를 두지 않습니다};

\end{tikzpicture}
\end{document}
